% §3 The Trajectory Grammar
We now define the paper's central formal object: a grammar over voice-agent conversations whose
runs are trajectories and whose terminal classes are banking outcomes.
The definition is deliberately machine-level: every state and transition corresponds to an event
that a production cascade already emits, so grammaticality can be checked from telemetry alone.
Existing agent benchmarks score conversations only at full-episode granularity
\citep{sierra2024taubench,webarena2023}, and workflow mining recovers structure from completed
logs \citep{salesforce2025workflows}; neither makes the conversation state first-class while the
conversation is still running.
A machine-level grammar also gives the operations team a shared vocabulary: when a reviewer asks
why a call was handed off, the answer is a walk through $G$, not an anecdote.
Finally, the grammar is the substrate on which the rest of the paper builds: Section~4 stops on it,
Section~5 measures it, and Section~6 acts on it.

\subsection{Conversation States and Typed Transitions}

\begin{definition}[Trajectory Grammar]
\label{def:grammar}
A \emph{trajectory grammar} is a triple $G = (Q, \Sigma, \delta)$ where the state set
$Q = Q_p \cup T$ partitions into progress states
$Q_p = \{q_{greet}, q_{auth}, q_{intent}, q_{tool}, q_{clarify}, q_{confirm}, q_{close}\}$
and terminal outcome classes
$T = \{\theta_{cont}, \theta_{aban}, \theta_{hand}, \theta_{esc}\}$
(contained, abandoned, handoff, compliance-escape).
The alphabet $\Sigma$ is a finite set of \emph{typed} transitions (Table~\ref{tab:transitions}),
and $\delta : Q \times \Sigma \rightharpoonup Q$ is a partial transition function.
Each transition type carries a direction label: \emph{progress} types advance the conversation
toward a service goal, and \emph{regression} types return the run to a previously visited
non-terminal state.
\end{definition}

\begin{figure}[t]
  \centering
  \resizebox{\linewidth}{!}{%
    % fig_grammar.tex — the trajectory grammar: states, typed transitions, terminal classes
\begin{tikzpicture}[
  node distance=1.15cm and 1.35cm,
  state/.style={circle, draw, minimum size=1.05cm, font=\small, fill=blue!5},
  terminal/.style={rectangle, rounded corners, draw, minimum width=1.5cm, minimum height=0.62cm, font=\small\itshape},
  every edge/.append style={-Stealth, font=\scriptsize},
  looplabel/.style={font=\scriptsize, align=center},
]
\node[state] (gr) {gr};
\node[state, right=of gr] (au) {au};
\node[state, right=of au] (in) {in};
\node[state, right=of in] (tl) {tl};
\node[state, right=of tl] (cl) {cf};
\node[state, right=of cl] (co) {cl};

\node[terminal, below=2.1cm of in] (cont) {contained};
\node[terminal, left=1.1cm of cont] (aban) {abandoned};
\node[terminal, right=1.1cm of cont] (hand) {handoff};
\node[terminal, below=0.7cm of cont] (comp) {compliance-escape};

\draw (gr) edge[bend left=18] node[above]{answer} (au);
\draw (au) edge[bend left=18] node[above]{verify} (in);
\draw (in) edge[bend left=18] node[above]{request} (tl);
\draw (tl) edge[bend left=18] node[above]{result} (cl);
\draw (cl) edge[bend left=18] node[above]{confirm} (co);
\draw (co) edge node[right]{close} (cont);

% regression (cycle) edges
\draw (tl) edge[bend left=30] node[looplabel, below]{tool-fail$\to$retry} (in);
\draw (cl) edge[bend left=30] node[looplabel, below]{mismatch$\to$clarify} (in);
\draw (au) edge[bend right=35, looseness=1.4] node[looplabel, below left]{re-auth} (au);
\draw (co) edge[bend left=22] node[looplabel, above, sloped]{regress} (tl);

% terminal edges
\draw (cl) edge node[right]{no-confirm} (aban);
\draw (cl) edge node[left]{defer} (hand);
\draw (tl) edge node[looplabel, right]{policy breach} (comp);

% drift annotations
\begin{scope}[on background layer]
  \node[fit=(in)(tl), draw, dashed, rounded corners, inner sep=6pt, label={[font=\scriptsize]below:drift zone: cycle depth $D$ counted here}] {};
\end{scope}
\end{tikzpicture}%
  }
  \caption{The trajectory grammar $G$.
           Progress states (rounded) form the service spine; typed transitions ( labelled
           arrows) carry the eight event classes of Table~\ref{tab:transitions};
           regression types open cycles back into earlier states;
           terminal classes (squares) absorb every well-formed run.}
  \label{fig:grammar}
\end{figure}

The seven progress states cover the lifecycle of a banking call: identification, authentication,
intent capture, tool execution, clarification, confirmation, and closing.
The terminal classes are not generic success labels but the four production KPIs of conversational
banking: containment rate, abandonment rate, handoff rate, and compliance-escape rate.
Each class has an operational definition in telephony terms.
$\theta_{cont}$ (contained) holds when the run reaches a terminal through $q_{close}$ after a
confirmed service goal.
$\theta_{aban}$ (abandoned) holds when the caller disconnects before any terminal confirmation,
which in production surfaces as a hang-up event.
$\theta_{hand}$ (handoff) holds when an \texttt{escalation} transition fires, transferring the
caller to a human teller.
$\theta_{esc}$ (compliance-escape) holds when the run reaches a terminal confirmation without the
guarded \texttt{disclosure} transition having fired, which is precisely the false-success pattern
reported for tool agents \citep{falsesuccess2026}.
Industrial practice measures these quantities post hoc from telephony telemetry
\citep{nubank2026,silentabandonment2024}; under the grammar they become properties of runs.
This matters because final labels themselves mislead: agents can assert success while the task
failed \citep{falsesuccess2026}, so a terminal class must be derivable from the event stream
rather than from the agent's own closing statement.

\begin{table}[t]
\centering
\caption{Typed transitions of the trajectory grammar.
         Emitters are production components; regression types increment cycle depth.}
\label{tab:transitions}
\small
\begin{tabular}{lllll}
\toprule
Type & Direction & Source $\to$ Target & Emitter & Class \\
\midrule
\texttt{user-answer}   & progress   & $q_{auth/intent/clarify/confirm} \to$ next spine state & orchestrator LM      & dialogue \\
\texttt{agent-clarify} & either     & $q_{intent/confirm} \leftrightarrow q_{clarify}$       & orchestrator LM      & dialogue \\
\texttt{tool-call}     & progress   & $q_{tool} \to q_{tool}$ (iteration)                    & orchestrator         & system \\
\texttt{tool-fail}     & regression & $q_{tool} \to q_{tool}$ (retry)                        & backend tool         & system \\
\texttt{endpoint-error}& regression & $q_{intent/confirm} \to q_{intent}$                    & ASR endpointer       & voice \\
\texttt{user-bargein}  & event      & any $q_p \to q_p$ (audio channel)                      & barge-in detector    & voice \\
\texttt{disclosure}    & progress   & $q_{confirm} \to q_{close}$ (guarded)                  & compliance layer     & compliance \\
\texttt{escalation}    & progress   & any $q_p \to \theta_{hand}$                            & policy layer         & dialogue \\
\bottomrule
\end{tabular}
\end{table}

Table~\ref{tab:transitions} assigns every transition type to the component that emits it, which
is what makes the grammar checkable from telemetry: no state requires human annotation to
identify.
Two transition types are voice-physics events rather than dialogue acts: \texttt{endpoint-error}
fires when the endpointer commits a turn boundary too early or too late, and \texttt{user-bargein}
fires when the caller talks over the agent.
Both are observable from the same detection stages that bound response onset and cessation in the
published floors \citep{garcia2026latencyfloor,garcia2026cessationfloor}.
The \texttt{agent-clarify} type is marked \emph{either} because a clarification is progress when
it resolves an ambiguity and a regression when it re-opens a settled state; the grammar does not
prejudge which, and the cycle count of Definition~\ref{def:cycledepth} is what decides.
The \texttt{user-bargein} type leaves the state unchanged because a barge-in is a channel event,
not a conversation state; its conversational meaning is an observation for the posterior of
Definition~\ref{def:posterior}, not a state change.
This separation keeps the grammar auditable: every state change has a named emitter in the
orchestration stack, and channel events never mutate conversation state by themselves.

\begin{definition}[Run]
\label{def:run}
A \emph{run} of $G$ is a finite alternating sequence
$\sigma = q_0\, t_1\, q_1\, \cdots\, t_n\, q_n$ with $q_0 = q_{greet}$, $q_n \in T$, and each
$(t_i, q_i)$ permitted by $\delta$.
A \emph{partial run} is any proper prefix $\sigma_{\leq j}$ of a well-formed run that ends in a
non-terminal state.
The \emph{true outcome} of $\sigma_{\leq j}$ is the terminal class of the well-formed run that
extends it, which is unique by absorption but unobserved while the conversation runs.
Early outcome prediction is exactly the problem of guessing this unobserved terminal class from
$\sigma_{\leq j}$, and Section~4 states when that guess may be acted on.
$\sigma$ is \emph{well-formed} if every terminal class is absorbing, i.e., no transition departs
from a state in $T$.
Well-formedness is a safety property of the orchestration policy, not of the data: a run that
continues after a terminal class is a policy bug, and the grammar makes such bugs detectable by
inspection of $\delta$.
\end{definition}

\subsection{Cycles and the Drift Statistic}

\begin{definition}[Cycle and cycle depth]
\label{def:cycledepth}
A \emph{cycle} of a run $\sigma$ is a maximal sub-run that departs from a state $q$ and later
returns to $q$, where $q$ is non-terminal.
The \emph{cycle depth} $D(\sigma)$ is the number of cycles in $\sigma$.
\end{definition}

The intuition is operational: one clarification after a failed balance lookup is normal service,
whereas a third return to $q_{clarify}$ is a conversation going in circles.
Cycle depth turns this intuition into a countable statistic that separates two failure modes that
conversation-length measures conflate.
Three properties fix the statistic's behaviour.
First, $D(\sigma)$ counts returns, not repetitions: iterating \texttt{tool-call} inside
$q_{tool}$ is spine progress, while returning to $q_{intent}$ after having confirmed is drift.
Second, $D$ is monotone along a run, so the monitor maintains it online at zero detection cost
beyond event counting.
Third, $D$ is defined identically for real and simulated runs, which makes it portable to the
protocols of \citet{sierra2024taubench} and the workflow corpora of
\citet{salesforce2025workflows} without re-labelling.

\begin{proposition}[Longer is not off-path]
\label{prop:longer_offpath}
There exist well-formed runs $\sigma_L$ and $\sigma_D$ of equal length $n$ with
$D(\sigma_L) = 0$ and $D(\sigma_D) \geq 1$.
Hence run length alone determines neither drift nor the terminal class, and any outcome detector
keyed to conversation duration confuses the two.
The distinction has direct operational value: a long call that is merely tool-heavy should not
page a supervisor, while a short call with two clarification cycles may already be failing.
\end{proposition}

\begin{proof}
Construct $\sigma_L$ as the direct spine
$q_{greet} \to q_{auth} \to q_{intent} \to q_{tool} \to q_{confirm} \to q_{close} \to \theta_{cont}$
padded with $n - 6$ \texttt{tool-call} iterations inside $q_{tool}$; no state is departed and
revisited, so $D(\sigma_L) = 0$.
Construct $\sigma_D$ with the same number of transitions but routing two of them through
$q_{confirm} \to q_{clarify} \to q_{confirm}$ via \texttt{agent-clarify}, removing an equal number
of \texttt{tool-call} iterations to hold length fixed; this sub-run departs from and returns to
$q_{confirm}$, so $D(\sigma_D) \geq 1$.
Both runs are well-formed and of length $n$, which proves the claim.
\end{proof}

\begin{figure}[t]
  \centering
  \resizebox{\linewidth}{!}{%
    % fig_cycles.tex — contained vs drifted trajectories as grammar walks (cycle depth)
\begin{tikzpicture}[xscale=1.0, yscale=1.0]
% Panel A: contained call (low cycle depth)
\begin{scope}
  \node[font=\small\bfseries, anchor=west] at (-0.4,2.3) {(a) contained: $D=0$};
  % turn dots along a line
  \foreach \i in {0,...,9} { \fill[blue!60!black] (\i,0) circle (1.7pt); }
  \draw[gray!60] (-0.4,0) -- (9.4,0);
  % state strip
  \node[font=\scriptsize, anchor=west] at (-0.4,-0.55) {gr $\to$ au $\to$ in $\to$ tl $\to$ cf $\to$ cl};
  \node[font=\scriptsize, align=left, anchor=west] at (-0.4,-1.25) {monotone progress; no repeated states; outcome decidable at $k_{\mathrm{dec}}=6$};
  \node[font=\scriptsize, blue!60!black, anchor=west] at (-0.4,0.45) {posterior crosses $1-\delta$ inside the tool loop};
\end{scope}
% Panel B: drifted call (high cycle depth)
\begin{scope}[yshift=-3.3cm]
  \node[font=\small\bfseries, anchor=west] at (-0.4,2.3) {(b) drifted: $D=3$};
  \draw[gray!60] (-0.4,0) -- (9.4,0);
  \foreach \i in {0,...,12} { \fill[red!60!black] (\i*0.72,0) circle (1.7pt); }
  % cycles drawn as arcs back
  \draw[-Stealth, red!60!black] (2.2,0.12) .. controls (3.0,0.9) and (4.2,0.9) .. (5.0,0.12)
    node[midway, above, font=\scriptsize] {retry};
  \draw[-Stealth, red!60!black] (4.3,-0.12) .. controls (5.1,-1.0) and (6.3,-1.0) .. (7.1,-0.12)
    node[midway, below, font=\scriptsize] {re-clarify};
  \draw[-Stealth, red!60!black] (5.8,0.12) .. controls (6.6,1.2) and (7.8,1.2) .. (8.6,0.12)
    node[midway, above, font=\scriptsize] {regress};
  \node[font=\scriptsize, align=left, anchor=west] at (-0.4,-1.75) {same turn count as 1.4$\times$ a contained call, but three returns to visited states: off-path};
  \node[font=\scriptsize, red!60!black, anchor=west] at (-0.4,-2.45) {cycle-depth statistic fires drift warning at $D\geq 2$, before the posterior alone would};
\end{scope}
\end{tikzpicture}%
  }
  \caption{Cycle depth as drift.
           Two runs of equal length: the contained run (top) iterates \texttt{tool-call} within
           $q_{tool}$ and stays on the spine; the drifted run (bottom) opens a
           $q_{confirm} \to q_{clarify} \to q_{confirm}$ cycle.
           Length is identical; only cycle depth distinguishes them.}
  \label{fig:cycles}
\end{figure}

Cycle depth would be unbounded without a turn-scale trigger, and the published latency program
supplies one.

\begin{lemma}[Cycle-length bound]
\label{lem:cyclebound}
Fix a per-response deadline $\tau_{resp}$ and the latency parameters of the cascade pipeline, and
let $\Lambda_k = \lambda_{asr} + k \cdot \lambda_{lm} + (k-1) \cdot \lambda_{tool} + \lambda_{tts}$
be the $k$-step silence gap of \citet{garcia2026latencyfloor}.
Let $k^* = \max\{k : \Lambda_k \leq \tau_{resp}\}$.
Then:
(i) any run segment between two consecutive user turns that emits more than $k^*$
\texttt{tool-call} transitions violates the response deadline, and is therefore off-path for every
deadline-conforming policy;
(ii) under the patience-hazard model of \citet{silentabandonment2024}, the probability that the
caller remains on the line decays with each deadline breach, so runs far beyond $k^*$ terminate
in $\theta_{aban}$ with high probability.
\end{lemma}

\begin{proof}
For (i), each \texttt{tool-call} iteration inside one orchestrated turn adds one LM step and one
tool round-trip to the user-perceived silence gap, which grows linearly in $k$ by the compounding
result of \citet{garcia2026latencyfloor}; emitting more than $k^*$ such transitions therefore
produces a silence gap exceeding $\tau_{resp}$, which no deadline-conforming policy permits.
For (ii), each breached deadline raises the instantaneous hang-up hazard estimated from
contact-center patience data \citep{silentabandonment2024}; iterating the hazard over repeated
breaches drives the survival probability to zero, which yields the high-probability claim.
\end{proof}

\begin{remark}
\label{rem:drift_trigger}
Lemma~\ref{lem:cyclebound} makes drift detection an event-counting rule, not a timeout: the
monitor raises a drift alarm at the $(k^*{+}1)$-th \texttt{tool-call} event, in real time and
without waiting for the response to complete.
Because \texttt{tool-call} events are emitted by the orchestrator itself, the count is exact; the
only probabilistic component is the caller-patience channel in claim (ii).
The alarm's own detection latency is the monitor's decision time, which Section~4 bounds by a
stopping rule and Section~5 accounts against the published floors.
Two properties make $k^*$ usable in production rather than merely in theory.
It is monotone in the deadline: relaxing $\tau_{resp}$ raises $k^*$, so the drift trigger adapts
to policy rather than being hard-coded.
And it is auditable: given the latency table of the deployed cascade, $k^*$ is a one-line
computation that a reviewer can re-derive from the floor arithmetic of
\citet{garcia2026latencyfloor} alone.
\end{remark}

\subsection{Multi-Channel Observation Model}

\begin{definition}[Observation]
\label{def:observation}
At turn $k$ the monitor observes the tuple
$o_k = (x_k, s_k, e_k, f_k)$:
the partial transcript $x_{1:k}$;
the turn-state sequence $s_{1:k}$ labelling each utterance as completion, backchannel, or
interruption \citep{x2turn2026};
the tool-event log $e_{1:k}$ emitted by the orchestrator; and
the acoustic flags $f_{1:k}$ recording endpointing errors and barge-ins
\citep{garcia2026cessationfloor}.
Every channel is streaming: each component updates incrementally, so $o_k$ is available within one
turn of real time and no channel requires the conversation to have ended.
\end{definition}

The channels are not interchangeable, and the transcript channel is the weakest one.
Text-only turn-end decisions sit near chance on Spanish hard cuts (accuracy $0.40$--$0.54$), as
measured in our semantic-VAD study \citep{garcia2026asking}; a grammar that conditioned only on
$x_{1:k}$ would therefore inherit that blindness at conversation scale.
The turn-state and tool-event channels carry exactly the information the text channel lacks:
$s_k$ distinguishes completed intents from backchannels, and $e_k$ reports system-side progress
that never surfaces in the transcript at all.

\begin{definition}[Outcome posterior]
\label{def:posterior}
Given the observation history $o_{1:k}$, the outcome posterior is
$p(\theta \mid o_{1:k})$ over terminal classes $\theta \in T$.
\end{definition}

The posterior is defined to be recomputable after every event, not only at turn boundaries,
because every channel updates incrementally in a streaming cascade.
Section~4 places a sequential stopping rule over $p(\theta \mid o_{1:k})$ and derives when the
monitor may act on it.
The price of early warning is paid per turn inside the same floor arithmetic as the response
itself: the published cloud decision latencies (p99 $178$--$268$~ms) \citep{garcia2026asking}
bound what each posterior update adds, and every intervention triggered by the posterior must
additionally clear the cessation floor before the agent's audio actually stops
\citep{garcia2026cessationfloor}.
